\section{Computer Networks}
\label{sec:qna-cn}

\subsection*{Models and layers}

\question{Explain the OSI model layers.}
\answer{Seven layers, bottom to top: \textbf{Physical} (bits on the wire),
\textbf{Data Link} (frames, MAC), \textbf{Network} (packets, IP, routing),
\textbf{Transport} (end-to-end, TCP/UDP), \textbf{Session} (dialog
control), \textbf{Presentation} (encoding/encryption), \textbf{Application}
(HTTP, DNS). Mnemonic: ``Please Do Not Throw Sausage Pizza Away.''}

\question{Difference between the OSI and TCP/IP models?}
\answer{OSI is a 7-layer reference model; TCP/IP is the 4-layer
(Link, Internet, Transport, Application) practical model the internet
actually uses. TCP/IP merges OSI's top three into Application and the
bottom two into Link. OSI is for teaching; TCP/IP is for building.}

\subsection*{Transport}

\question{TCP vs UDP?}
\answer{TCP is connection-oriented, reliable, ordered, and flow/congestion
controlled -- at the cost of overhead and latency (used for web, email,
file transfer). UDP is connectionless, unreliable, and unordered but fast
and lightweight (used for video, VoIP, gaming, DNS). Choose by whether you
need reliability or low latency.}

\question{Explain the TCP three-way handshake.}
\answer{To establish a connection: the client sends \texttt{SYN}, the
server replies \texttt{SYN-ACK}, and the client sends \texttt{ACK}. This
synchronises sequence numbers in both directions and confirms both sides
can send and receive before data flows.}

\question{How does TCP ensure reliability?}
\answer{Through sequence numbers (ordering), acknowledgements and
retransmission (lost-packet recovery), checksums (corruption detection),
flow control via a sliding window (don't overwhelm the receiver), and
congestion control (don't overwhelm the network).}

\question{What is the difference between flow control and congestion
control?}
\answer{Flow control prevents a fast sender from overrunning a slow
\emph{receiver}, using the receiver's advertised window. Congestion control
prevents overloading the \emph{network} itself, using algorithms like slow
start and AIMD (additive increase, multiplicative decrease).}

\subsection*{Application and addressing}

\question{What happens when you type a URL and press Enter?}
\answer{DNS resolves the domain to an IP; the browser opens a TCP
connection (and a TLS handshake for HTTPS); it sends an HTTP request; the
server responds; the browser parses HTML and fetches further resources
(CSS, JS, images), then renders. Caching, redirects, and keep-alive
connections optimise the steps.}

\question{What is DNS and how does resolution work?}
\answer{DNS maps human-readable domain names to IP addresses. Resolution is
hierarchical: the resolver queries a root server, then the TLD server
(e.g.\ \texttt{.com}), then the domain's authoritative server, caching
results along the way to speed future lookups.}

\question{Difference between HTTP and HTTPS?}
\answer{HTTPS is HTTP over TLS: it encrypts the traffic, authenticates the
server via certificates, and protects integrity, preventing eavesdropping
and tampering. HTTP is plaintext. HTTPS adds a TLS handshake but is now the
default for the web.}

\question{TCP vs UDP -- which does HTTP use, and is that changing?}
\answer{HTTP/1.1 and HTTP/2 run over TCP. HTTP/3 runs over QUIC, which is
built on \emph{UDP}, to avoid TCP's head-of-line blocking and speed up
connection setup -- an example of choosing UDP plus custom reliability for
performance.}

\subsection*{Devices and security}

\question{Difference between a router, a switch, and a hub?}
\answer{A hub broadcasts to all ports (Layer 1, dumb). A switch forwards
frames only to the destination port using MAC addresses (Layer 2). A router
connects different networks and forwards packets by IP address, choosing
routes (Layer 3).}

\question{What is NAT?}
\answer{Network Address Translation maps many private IP addresses to one
(or few) public IPs at the router, conserving scarce IPv4 addresses and
adding a layer of isolation. The router rewrites addresses/ports and tracks
the mapping to route replies back.}

\question{What is the difference between symmetric and asymmetric
encryption?}
\answer{Symmetric uses one shared secret key for both encryption and
decryption -- fast, but key distribution is hard. Asymmetric uses a
public/private key pair -- solving key distribution but slower. TLS combines
them: asymmetric to exchange a session key, then symmetric for bulk data.}
